% ============================================================================
%  Truncated & Renormalized 3D ETAS Model — Mathematical Description
%  ----------------------------------------------------------------------------
%  Compile with any standard LaTeX distribution:
%      pdflatex MODEL.tex   (run twice for cross-references)
%  or  xelatex  MODEL.tex
%
%  This document mirrors MODEL.md and adds full closed-form proofs of the
%  renormalization constants and their total derivatives.
%
%  Sources verified against:
%    [GZZ15]  Guo, Zhuang & Zhou (2015), Geophys. J. Int. 203, 366-372
%             (ggv319.pdf)   -- hypocentral (3D) depth kernel
%    [J19]    Jalilian (2019), J. Stat. Software 88, Code Snippet 1
%             (Jalilian.pdf) -- space-time ETAS R package
% ============================================================================
\documentclass[11pt,a4paper]{article}

% ---- Math & symbols --------------------------------------------------------
\usepackage{amsmath,amssymb,amsthm,mathtools}
\usepackage{bm}
\usepackage{booktabs}
\usepackage{array}
\usepackage{microtype}
\usepackage[hmargin=2.4cm,vmargin=2.5cm]{geometry}
\usepackage{caption}
\usepackage{enumitem}
\setlist{itemsep=2pt,topsep=2pt}

% ---- Hyperlinks / code listings --------------------------------------------
\usepackage[round]{natbib}
\usepackage{hyperref}
\hypersetup{colorlinks=true,linkcolor=blue!55!black,
            urlcolor=blue!55!black,citecolor=blue!55!black}
\usepackage{listings}
\lstset{basicstyle=\ttfamily\footnotesize,breaklines=true,columns=fullflexible,
        frame=single,framerule=0.3pt,rulecolor=\color{gray!50},
        xleftmargin=4pt,xrightmargin=4pt}

% ---- Theorem-like environments ---------------------------------------------
\theoremstyle{plain}
\newtheorem{proposition}{Proposition}
\newtheorem{lemma}{Lemma}
\theoremstyle{definition}
\newtheorem{defn}{Definition}

% ---- Convenience macros ----------------------------------------------------
\newcommand{\dt}{\delta t}
\newcommand{\eps}{\varepsilon}
\newcommand{\Gnorm}{G_{\mathrm{norm}}}
\newcommand{\Fnorm}{F_{\mathrm{norm}}}
\newcommand{\Hnorm}{H_{\mathrm{norm}}}
\newcommand{\Hint}{H_{\mathrm{int}}}
\newcommand{\Gint}{G_{\mathrm{int}}}
\newcommand{\Sint}{S_{\mathrm{int}}}
\newcommand{\Tmax}{T_{\max}}
\newcommand{\Rmax}{R_{\max}}
\newcommand{\Zmax}{Z_{\max}}
\newcommand{\sigj}{\sigma_j}
\newcommand{\bD}{\mathcal{H}}
\newcommand{\R}{\mathbb{R}}
\newcommand{\E}{\mathbb{E}}
\newcommand{\dd}{\mathrm{d}}

% ---- Title -----------------------------------------------------------------
\title{\textbf{Truncated \& Renormalized 3D ETAS Model}\\
       \large Mathematical Description and Implementation Notes}
\author{ETAS Python}
\date{}

\begin{document}
\maketitle

\begin{abstract}
This document gives the mathematical specification of the \emph{Truncated \&
Renormalized} Epidemic-Type Aftershock Sequence (ETAS) model implemented in the
\texttt{etas\_python} package. The base space-time model follows \citet{O98}
and the ETAS R package of Jalilian~\citep{J19}; the hypocentral depth extension
follows \citet{GZZ15}. We give full closed-form derivations
of the truncation cutoffs, the renormalization constants, and their
\emph{total} derivatives through the cutoff constraint, which the maximum
likelihood optimizer requires. Throughout, the standard (un-truncated) model is
recovered exactly as the truncation threshold $\eps\to 0$.
\end{abstract}

\tableofcontents

% ============================================================================
\section{Model Overview}
% ============================================================================

The ETAS model describes earthquake occurrence as a self-exciting
spatio-temporal point process. The conditional intensity at location $(x,y,z)$
and time $t$ is
\begin{equation}\label{eq:intensity}
\lambda(t,x,y,z\mid\bD_t)
= \mu(x,y,z)
  + \sum_{i:\,t_i<t} \kappa(m_i)\, g(\dt_{ij};c,p)\,
        f(r_{ij}; m_i,D,q,\gamma)\, h(z_j; z_i,\eta,\Zmax),
\end{equation}
where $\bD_t$ is the history up to time $t$, $\mu$ is the background rate,
$\kappa(m_i)=A e^{\alpha m_i}$ is the productivity, $g$ the Omori--Utsu temporal
kernel, $f$ the spatial kernel (power-law or Gaussian), and $h$ the optional
3D depth kernel (Beta density on $[0,\Zmax]$). This form is exactly eqs.~(2)
of Jalilian~\cite{J19} and eqs.~(3)--(4) of Guo \emph{et al.}~\cite{GZZ15},
augmented by the background term $\mu(x,y,z)$.

\paragraph{Magnitude convention.}
Both source papers parameterize the kernels with the offset
$m-m_0$ (e.g.\ $\kappa(m)=A e^{\alpha(m-m_0)}$, $\sigma_j=D e^{\gamma(m_j-m_0)}$).
In this implementation the magnitude offset $m_0$ is subtracted once, at
catalog construction time, so the array $m$ already stores
$m_i-m_0$. Hence every formula below is written in terms of the offset
magnitude and the explicit ``$-m_0$'' is omitted. This is a notational
convention only; the mathematics is unchanged.

\paragraph{Document scope.}
This document describes the \emph{Truncated \& Renormalized} extension that
makes the model scalable to large catalogues ($N>10^5$) while remaining a
valid probability model. All closed forms below are derived in
Sections~\ref{sec:kernels}--\ref{sec:renorm-grad} and verified against the
reference implementation (\texttt{etas/src/renorm.py},
\texttt{etas/src/lambda\_funcs.py}).

% ============================================================================
\section{Kernel Functions}\label{sec:kernels}
% ============================================================================

% ----------------------------------------------------------------------------
\subsection{Temporal Omori--Utsu Kernel}\label{sec:omori}

\begin{equation}\label{eq:omori}
g(\dt; c,p) = \frac{p-1}{c}\Bigl(1+\frac{\dt}{c}\Bigr)^{-p},
\qquad \dt\ge 0,\; c>0,\; p>1.
\end{equation}
This is the modified Omori law used by both \cite{J19} and \cite{GZZ15}.

\begin{proposition}[Normalization]\label{prop:omori-norm}
$g$ of \eqref{eq:omori} integrates to $1$ over $[0,\infty)$.
\end{proposition}
\begin{proof}
With $\tau=\dt/c$ and $a=1-p<0$,
$\frac{\dd}{\dd\tau}(1+\tau)^a = a(1+\tau)^{a-1}$, so
\[
\int_0^\infty g(\dt)\,\dd\dt
= \frac{p-1}{c}\int_0^\infty\Bigl(1+\frac{\dt}{c}\Bigr)^{-p}\dd\dt
= (p-1)\int_0^\infty(1+\tau)^{-p}\dd\tau
= (p-1)\Bigl[\tfrac{(1+\tau)^{1-p}}{1-p}\Bigr]_0^\infty
= 1. \qedhere
\]
\end{proof}

% ----------------------------------------------------------------------------
\subsection{Spatial Power-Law Kernel \texttt{(mver=1)}}\label{sec:pl-kernel}

\begin{equation}\label{eq:pl-kernel}
f(r; m_j,D,\gamma,q)
= \frac{q-1}{\pi\sigj}\Bigl(1+\frac{r^2}{\sigj}\Bigr)^{-q},
\qquad \sigj = D e^{\gamma m_j},\; q>1.
\end{equation}
The angular integral over $\R^2$ gives $2\pi r\,\dd r$, hence
\begin{equation}\label{eq:pl-norm}
\int_{\R^2} f\,\dd x\dd y
= \int_0^\infty \frac{2(q-1)\,r/\sigj}{(1+r^2/\sigj)^q}\,\dd r
= \Bigl[-\bigl(1+r^2/\sigj\bigr)^{1-q}\Bigr]_0^\infty = 1.
\end{equation}

% ----------------------------------------------------------------------------
\subsection{Gaussian Spatial Kernel \texttt{(mver=2)}}\label{sec:gauss-kernel}

\begin{equation}\label{eq:gauss-kernel}
f(r; m_j,D,\gamma)
= \frac{1}{2\pi\sigj}\exp\Bigl(-\frac{r^2}{2\sigj}\Bigr),
\qquad \sigj = D e^{\gamma m_j}.
\end{equation}

% ----------------------------------------------------------------------------
\subsection{3D Depth Kernel (Beta density on $[0,\Zmax]$)}\label{sec:depth-kernel}

Define normalized depths $u=z/\Zmax$, $v=z'/\Zmax$. The kernel is
\begin{equation}\label{eq:depth-kernel}
h(u;v,\eta) =
\frac{u^{\eta v}(1-u)^{\eta(1-v)}}
     {\Zmax\, B\bigl(\eta v+1,\;\eta(1-v)+1\bigr)},
\end{equation}
where $B(a,b)=\Gamma(a)\Gamma(b)/\Gamma(a+b)$ is the Beta function. This is
eq.~(5) of \cite{GZZ15}.

\begin{proposition}[Normalization]\label{prop:depth-norm}
For every parent depth $v\in[0,1]$, $\int_0^{\Zmax} h(z;z')\,\dd z = 1$.
\end{proposition}
\begin{proof}
With $u=z/\Zmax$, $\dd z=\Zmax\,\dd u$, and $a=\eta v+1$, $b=\eta(1-v)+1$,
\[
\int_0^{\Zmax} h\,\dd z
= \frac{1}{B(a,b)}\int_0^1 u^{a-1}(1-u)^{b-1}\,\dd u
= \frac{B(a,b)}{B(a,b)} = 1. \qedhere
\]
The parameter $\eta\ge 0$ controls concentration: $\eta=0$ gives a uniform
depth distribution, large $\eta$ concentrates offspring near the parent depth.
\end{proof}

% ============================================================================
\section{Truncation and Analytic Renormalization}\label{sec:trunc}
% ============================================================================

For large catalogues the double-sum over $(i,j)$ dominates as $O(N^2)$. We
truncate each kernel at a threshold $\eps>0$ where the kernel value is
negligible, and \emph{renormalize} so the truncated kernel remains a valid
probability density on its surviving support. All constants below are
closed-form (no numerical integration), so they cost one vectorized evaluation
per likelihood call.

% ----------------------------------------------------------------------------
\subsection{Temporal Truncation}\label{sec:trunc-t}

\begin{lemma}[Temporal cutoff and normalizer]\label{lem:t-norm}
Solving $g(\Tmax)=\eps_t$ for $\Tmax$ gives
\begin{equation}\label{eq:tmax}
\Tmax = c\Bigl[\Bigl(\frac{p-1}{c\eps_t}\Bigr)^{1/p}-1\Bigr],
\end{equation}
and the integral over $[0,\Tmax]$ is
\begin{equation}\label{eq:gnorm}
\Gnorm(c,p,\eps_t) = \int_0^{\Tmax} g(\tau)\,\dd\tau
= 1-\Bigl(1+\frac{\Tmax}{c}\Bigr)^{1-p} = 1-u^{1-p},
\end{equation}
where $u=\bigl((p-1)/(c\eps_t)\bigr)^{1/p}$.
\end{lemma}
\begin{proof}
$g(\Tmax)=\eps_t$ means $(p-1)/c\cdot(1+\Tmax/c)^{-p}=\eps_t$, hence
$1+\Tmax/c=\bigl((p-1)/(c\eps_t)\bigr)^{1/p}$, which is \eqref{eq:tmax}. The
integral follows by Proposition~\ref{prop:omori-norm} applied to the upper
limit $\Tmax$:
$\int_0^{\Tmax} g = (p-1)\bigl[\tfrac{(1+\tau)^{1-p}}{1-p}\bigr]_0^{\Tmax/c}
= 1-(1+\Tmax/c)^{1-p}$.
\end{proof}

% ----------------------------------------------------------------------------
\subsection{Spatial Truncation}\label{sec:trunc-s}

\begin{lemma}[Spatial cutoff and normalizer]\label{lem:s-norm}
For the power-law kernel, solving $f(\Rmax(m_j)\mid m_j)=\eps_s$ gives
\begin{equation}\label{eq:rmax}
\Rmax(m_j) = \sqrt{\sigj\Bigl[\Bigl(\frac{q-1}{\pi\sigj\eps_s}\Bigr)^{1/q}-1\Bigr]},
\end{equation}
and the disk integral is
\begin{equation}\label{eq:fnorm}
\Fnorm(m_j;D,\gamma,q,\eps_s) = \int_0^{\Rmax} 2\pi r\, f(r\mid m_j)\,\dd r
= 1-\Bigl(1+\frac{\Rmax^2}{\sigj}\Bigr)^{1-q} = 1-v^{1-q},
\end{equation}
where $v=\bigl((q-1)/(\pi\sigj\eps_s)\bigr)^{1/q}$.
\end{lemma}
\begin{proof}
From \eqref{eq:pl-kernel}, $f(\Rmax)=\eps_s$ gives
$(q-1)/(\pi\sigj)\cdot(1+\Rmax^2/\sigj)^{-q}=\eps_s$, hence
$1+\Rmax^2/\sigj=\bigl((q-1)/(\pi\sigj\eps_s)\bigr)^{1/q}$, which yields
\eqref{eq:rmax}. The disk integral is the finite-support version of
\eqref{eq:pl-norm}:
$\int_0^{\Rmax}2\pi r f\,\dd r
=\bigl[-(1+r^2/\sigj)^{1-q}\bigr]_0^{\Rmax}
=1-(1+\Rmax^2/\sigj)^{1-q}$.
\end{proof}

\begin{remark}\label{rem:gauss-no-fnorm}
The Gaussian kernel \eqref{eq:gauss-kernel} is integrated analytically inside
the polygon integral (Section~\ref{sec:space-time-int}), so no per-event
$\Fnorm$ is computed for \texttt{mver=2}; $\Fnorm\equiv 1$ on that path.
\end{remark}

% ----------------------------------------------------------------------------
\subsection{Renormalized Kernels}\label{sec:renorm-kernels}

\begin{equation}\label{eq:renorm-kernels}
\tilde g(\dt)=\frac{g(\dt)}{\Gnorm},\qquad
\tilde f(r\mid m_j)=\frac{f(r\mid m_j)}{\Fnorm(m_j)},
\end{equation}
which by Lemmas~\ref{lem:t-norm}--\ref{lem:s-norm} satisfy
$\int_0^{\Tmax}\tilde g=1$ and $\int_{\mathrm{disk}}\tilde f=1$.

\begin{proposition}[Recovery of the un-truncated model]\label{prop:limit}
As $\eps_t,\eps_s\to 0$, the cutoffs $\Tmax,\Rmax\to\infty$ and
$\Gnorm,\Fnorm\to 1$, so the truncated model collapses exactly onto the
standard (un-pruned) ETAS model.
\end{proposition}
\begin{proof}
As $\eps_t\to 0$, $(p-1)/(c\eps_t)\to\infty$, so $\Tmax\to\infty$ and
$u^{1-p}\to 0$, giving $\Gnorm\to 1$. The spatial case is identical with
$v^{1-q}\to 0$. Thus \eqref{eq:renorm-kernels} tends to $g$ and $f$.
\end{proof}

% ----------------------------------------------------------------------------
\subsection{Renormalized Conditional Intensity}\label{sec:renorm-intensity}

With renormalization, the intensity at event $j$ becomes
\begin{equation}\label{eq:intensity-j}
\lambda_j = \mu b_j
+ \sum_{\substack{i<j \\ \dt_{ij}\le\Tmax \\ r_{ij}\le\Rmax(m_i)}}
  \kappa(m_i)\,\frac{g(\dt_{ij})}{\Gnorm}\,
  \frac{f(r_{ij}\mid m_i)}{\Fnorm(m_i)}\, h(z_j\mid z_i),
\end{equation}
where the sum runs only over parents $i$ surviving \emph{both} cutoffs
(enforced by KDTree lookups; see Section~\ref{sec:kdtree}). Note that
$\Gnorm$ is a scalar while $\Fnorm$ is a vector indexed by parent magnitude
$m_i$.

% ----------------------------------------------------------------------------
\subsection{Space-Time Integral under Renormalization}\label{sec:space-time-int}

The expected number of offspring triggered by parent $j$ inside the study
window is
\begin{equation}\label{eq:lambda-j-int}
\Lambda_j = \kappa(m_j)\cdot
\frac{\Gint}{\Gnorm}\cdot
\frac{\Sint(m_j)}{\Fnorm(m_j)}\cdot
\frac{\Hint}{\Hnorm},
\end{equation}
where $\Gint$ is the analytic Omori--Utsu CDF increment over the temporal
study window (closed form below), $\Sint$ the polygon integral of $f$
(numerical, radial partitioning of \citet{O98} as in \citet{J19}), and
$\Hint/\Hnorm\equiv 1$ for the depth kernel by
Proposition~\ref{prop:depth-norm}. The temporal increment is
\begin{equation}\label{eq:gi}
\Gint =
\begin{cases}
1-\bigl(1+(T_{\mathrm{end}}-t_j)/c\bigr)^{1-p}, & t_j \ge T_{\mathrm{start}},\\[4pt]
\bigl(1+(T_{\mathrm{start}}-t_j)/c\bigr)^{1-p}
-\bigl(1+(T_{\mathrm{end}}-t_j)/c\bigr)^{1-p}, & t_j < T_{\mathrm{start}},
\end{cases}
\end{equation}
which is exactly the closed form in \cite[\S3]{J19} (between eqs.~(5)--(6)).

% ============================================================================
\section{Gradients for Maximum-Likelihood Estimation}\label{sec:gradients}
% ============================================================================

All parameters are optimized in the \emph{sqrt-parameterisation}
\begin{equation}\label{eq:sqrt-param}
\phi_k = \theta_k^2,\qquad
\frac{\partial\phi_k}{\partial\theta_k}=2\theta_k,
\end{equation}
which enforces $\phi_k\ge 0$ automatically (no box constraints).

% ----------------------------------------------------------------------------
\subsection{Why the gradients are \emph{total}, not partial}\label{sec:total-vs-partial}

The likelihood recomputes $\Tmax$ and $\Rmax$ after each parameter perturbation
(they are functions of $c,p,D,\gamma,q$ through the cutoff constraint
$g(\Tmax)=\eps_t$, $f(\Rmax)=\eps_s$). Hence the gradient seen by the optimizer
is the \emph{total} derivative of $\Gnorm,\Fnorm$ through the cutoff, not the
partial derivative of $1-(1+\Tmax/c)^{1-p}$ holding $\Tmax$ fixed. To make this
dependence explicit we rewrite the constants in the cutoff-constraint
variables $u,v$ (Lemmas~\ref{lem:t-norm}--\ref{lem:s-norm}) and differentiate
directly.

% ----------------------------------------------------------------------------
\subsection{Gradients of the Renormalization Constants}\label{sec:renorm-grad}

\begin{proposition}[Temporal total derivatives]\label{prop:t-grad}
With $u=\bigl((p-1)/(c\eps_t)\bigr)^{1/p}$ and $\Gnorm=1-u^{1-p}$,
\begin{align}
\ln u &= \tfrac{1}{p}\ln\tfrac{p-1}{c\eps_t},\\
\frac{\dd u}{\dd c} &= -\frac{u}{pc},\qquad
\frac{\dd u}{\dd p} = u\Bigl[\frac{1}{p(p-1)}-\frac{\ln u}{p}\Bigr],\\
\frac{\dd\Gnorm}{\dd c} &= -\frac{p-1}{pc}\,u^{1-p},\label{eq:dGdc}\\
\frac{\dd\Gnorm}{\dd p} &= u^{1-p}\ln u-(1-p)\,u^{-p}\frac{\dd u}{\dd p}.\label{eq:dGdp}
\end{align}
\end{proposition}
\begin{proof}
From $\ln u=\tfrac{1}{p}[\ln(p-1)-\ln c-\ln\eps_t]$,
$\partial(\ln u)/\partial c=-1/(pc)$, hence $\dd u/\dd c=u\cdot(-1/(pc))$.
Also $\partial(\ln u)/\partial p
= -(\ln u)/p + 1/(p(p-1))$, hence $\dd u/\dd p=u\bigl[1/(p(p-1))-(\ln u)/p\bigr]$.
For $\Gnorm$, $\dd\Gnorm/\dd c=-(1-p)u^{-p}\dd u/\dd c=(p-1)u^{-p}\cdot u/(pc)
=-(p-1)u^{1-p}/(pc)$, giving \eqref{eq:dGdc}; \eqref{eq:dGdp} is the same
chain rule applied through $\dd u/\dd p$.
\end{proof}

\begin{proposition}[Spatial total derivatives]\label{prop:s-grad}
With $v=\bigl((q-1)/(\pi\sigma\eps_s)\bigr)^{1/q}$,
$\sigma=D e^{\gamma m}$, and $\Fnorm=1-v^{1-q}$,
\begin{align}
\ln v &= \tfrac{1}{q}\Bigl[\ln(q-1)-\ln(\pi\eps_s)-\ln\sigma\Bigr],\\
\frac{\dd v}{\dd D} &= -\frac{v}{qD},\qquad
\frac{\dd v}{\dd\gamma} = -\frac{vm}{q},\qquad
\frac{\dd v}{\dd q} = v\Bigl[\frac{1}{q(q-1)}-\frac{\ln v}{q}\Bigr],\\
\frac{\dd\Fnorm}{\dd D} &= -\frac{q-1}{qD}\,v^{1-q},\label{eq:dFdD}\\
\frac{\dd\Fnorm}{\dd\gamma} &= -\frac{q-1}{q}\,m\,v^{1-q},\label{eq:dFdg}\\
\frac{\dd\Fnorm}{\dd q} &= -\Bigl[(1-q)\,v^{-q}\frac{\dd v}{\dd q}-v^{1-q}\ln v\Bigr].\label{eq:dFdq}
\end{align}
\end{proposition}
\begin{proof}
Identical structure to Proposition~\ref{prop:t-grad}. Since
$\ln\sigma=\ln D+\gamma m$, $\partial(\ln\sigma)/\partial D=1/D$ and
$\partial(\ln\sigma)/\partial\gamma=m$, giving $\dd v/\dd D=-v/(qD)$ and
$\dd v/\dd\gamma=-vm/q$. The derivatives \eqref{eq:dFdD}--\eqref{eq:dFdq} are
the chain rule $\dd\Fnorm/\dd\phi=(q-1)v^{-q}\dd v/\dd\phi$ (using
$1-q=-(q-1)$), with the $\dd v/\dd q$ term left in implicit form because
$\ln v$ itself depends on $q$.
\end{proof}

% ----------------------------------------------------------------------------
\subsection{Intensity Gradient (Power-Law, \texttt{mver=1})}\label{sec:intensity-grad}

The contribution from parent $i$ to $\lambda_j$ is
\begin{equation}\label{eq:Iij}
I_{ij}
= \underbrace{A e^{\alpha m_i}}_{\kappa_i}\;
  \underbrace{\frac{(p-1)/c}{\Gnorm}\Bigl(1+\tfrac{\dt}{c}\Bigr)^{-p}}_{\tilde g}\;
  \underbrace{\frac{(q-1)/(\pi\sigj)}{\Fnorm(m_i)}\Bigl(1+\tfrac{r^2}{\sigj}\Bigr)^{-q}}_{\tilde f}\;
  h_{ij}.
\end{equation}
Let $k_i=A e^{\alpha m_i}$. The gradients w.r.t.\ the natural parameters are
\begin{align}
\frac{\partial I}{\partial A}     &= \frac{k_i}{A}\,\frac{g}{\Gnorm}\,\frac{f}{\Fnorm_i}\,h,\\
\frac{\partial I}{\partial c}     &= k_i\Bigl[\frac{1}{\Gnorm}\frac{\partial g}{\partial c}
                                     -\frac{g}{\Gnorm^2}\frac{\dd\Gnorm}{\dd c}\Bigr]\frac{f}{\Fnorm_i}\,h,\\
\frac{\partial I}{\partial p}     &= k_i\Bigl[\frac{1}{\Gnorm}\frac{\partial g}{\partial p}
                                     -\frac{g}{\Gnorm^2}\frac{\dd\Gnorm}{\dd p}\Bigr]\frac{f}{\Fnorm_i}\,h,\\
\frac{\partial I}{\partial D}     &= k_i\frac{g}{\Gnorm}
                                    \Bigl[\frac{1}{\Fnorm_i}\frac{\partial f}{\partial D}
                                          -\frac{f}{\Fnorm_i^2}\frac{\dd\Fnorm_i}{\dd D}\Bigr]h,\\
\frac{\partial I}{\partial\eta}   &= k_i\frac{g}{\Gnorm}\frac{f}{\Fnorm_i}\frac{\partial h}{\partial\eta},
\end{align}
with analogous $q,\gamma$ terms (replace the $D$ bracket by the
$\Fnorm_i$-quotient-rule form of Proposition~\ref{prop:s-grad}). The kernel
partials are
\begin{align}
\frac{\partial g}{\partial c} &= g\Bigl[-\frac{1}{c}+\frac{p\,\dt}{c(c+\dt)}\Bigr],
& \frac{\partial g}{\partial p} &= g\Bigl[\frac{1}{p-1}-\ln\Bigl(1+\tfrac{\dt}{c}\Bigr)\Bigr],\\
\frac{\partial f}{\partial D} &= \frac{f}{D}\Bigl[-1+q\Bigl(1-\tfrac{1}{1+r^2/\sigj}\Bigr)\Bigr],
& \frac{\partial f}{\partial\gamma} &= f\Bigl[-m+qm\Bigl(1-\tfrac{1}{1+r^2/\sigj}\Bigr)\Bigr].
\end{align}
The chain rule to the sqrt-parameters \eqref{eq:sqrt-param} multiplies each
term by $2\theta_k$ as the final step.

\begin{remark}[Implementation note]
A previous version of the code applied $k_i\cdot A\cdot e^{\alpha m_i}$ (a
double $e^{\alpha m_i}$) to the $c,p,D,q,\gamma,\eta$ derivatives. The
current implementation uses $k_i=A e^{\alpha m_i}$ once; this is verified by
the finite-difference tests in \texttt{tests/test\_gradient.py}.
\end{remark}

% ----------------------------------------------------------------------------
\subsection{Depth Kernel Gradient}\label{sec:depth-grad}

With $h(u;v)$ from \eqref{eq:depth-kernel},
\begin{equation}\label{eq:depth-grad}
\frac{\partial\ln h}{\partial\eta}
= v\ln u+(1-v)\ln(1-u)
  -\bigl[v\psi(\eta v+1)+(1-v)\psi(\eta(1-v)+1)-\psi(\eta+2)\bigr],
\end{equation}
where $\psi(x)=\dd\ln\Gamma(x)/\dd x$ is the digamma function, and
$\partial h/\partial\eta = h\cdot\partial\ln h/\partial\eta$.
The $\psi(\eta+2)$ term arises because the Beta normalization
$B(\eta v+1,\eta(1-v)+1)=\Gamma(\eta v+1)\Gamma(\eta(1-v)+1)/\Gamma(\eta+2)$
depends on $\eta$ through all three Gamma arguments.

% ============================================================================
\section{KDTree-Based Neighbour Pruning}\label{sec:kdtree}
% ============================================================================

The brute-force mask path ($O(N^2)$) is replaced by two lightweight lookups:

\begin{center}
\begin{tabular}{@{}llll@{}}
\toprule
Dimension & Method & Complexity (per query) & Implementation\\
\midrule
Temporal & \texttt{np.searchsorted} & $O(\log N)$ & Contiguous slice $[j-k,j)$\\
Spatial  & \texttt{cKDTree.query\_ball\_point} & $O(\log N+k)$ & KDTree on $(x,y)$\\
\bottomrule
\end{tabular}
\end{center}

\subsection{Class: \texttt{NeighborIndex}}

The class \texttt{etas.src.neighbors.NeighborIndex} is built once per EM
iteration:
\begin{enumerate}
\item \textbf{\texttt{\_\_init\_\_(x,y,t)}} --- stores sorted $t$ and builds a
      \texttt{cKDTree} on $(x,y)$.
\item \textbf{\texttt{set\_cutoffs(tau\_cut, r\_cut)}} --- configures the time
      window and spatial radius.
\item \textbf{\texttt{query(j,...)}} --- returns an \texttt{np.ndarray} of
      parent indices $i<j$ surviving both cutoffs.
\item \textbf{\texttt{query\_all(...)}} --- returns list-of-arrays for all $j$
      (used by \texttt{decluster.py}).
\end{enumerate}
The temporal query is a contiguous slice because the catalog is time-sorted;
the spatial query is a KDTree ball of radius $\Rmax(m_i)$, which depends on
the parent magnitude.

\subsection{Integration with the Likelihood}

When \texttt{nbr\_idx} is supplied to \texttt{lambda\_j}/\texttt{lambda\_j\_grad},
the parent loop iterates only over pre-pruned indices. Renorm constants are
precomputed once in \texttt{optimizer.\_precompute} and threaded through. When
no cutoffs are set the KDTree returns all $i<j$ and all norms are $1$, so by
Proposition~\ref{prop:limit} the model collapses exactly to the un-truncated
reference.

% ============================================================================
\section{EM Algorithm}\label{sec:em}
% ============================================================================

The model is fit by alternating:

\subsection{E-step: Stochastic Declustering (\texttt{decluster.py})}

For each event $j$, compute the background probability
\begin{equation}\label{eq:rho}
\rho_j = \frac{\mu b_j}{\lambda_j},
\end{equation}
assigning event $j$ to background with probability $\rho_j$ and to parent $i$
with probability $I_{ij}/\lambda_j$ (eqs.~(3)--(4) of \cite{J19}). The
``stochastic'' variant samples the parent once per event rather than storing
the full $N\times N$ responsibility matrix.

\subsection{M-step: DFP Quasi-Newton (\texttt{optimizer.py})}

The log-likelihood is
\begin{equation}\label{eq:loglik}
\ell(\theta) = \sum_{j=1}^{N}\ln\lambda_j(\theta)
              - \int_{\mathrm{window}}\lambda\,\dd t\,\dd x\,\dd y,
\end{equation}
and its gradient are computed using renormalized kernels and KDTree-pruned
parent lists (one shared \texttt{NeighborIndex} per EM iteration). The
Davidon--Fletcher--Powell (DFP) method updates $\theta$ with an approximate
inverse Hessian $H_k\to(\nabla^2\ell)^{-1}$ (Algorithm~1 of \cite{J19}).
The background rate is updated in closed form:
\begin{equation}\label{eq:mu-update}
\mu_{\mathrm{new}} = \frac{1}{N}\sum_{j=1}^{N}\rho_j.
\end{equation}

% ============================================================================
\section{Sqrt-Parameterisation}\label{sec:sqrt}
% ============================================================================

All ETAS parameters are non-negative. Rather than box constraints (slower for
quasi-Newton), we parameterise $\phi_k=\theta_k^2$ and optimize over
$\theta_k\in\R$. The chain rule
\begin{equation}\label{eq:chain}
\frac{\partial\ell}{\partial\theta_k}
= \frac{\partial\ell}{\partial\phi_k}\cdot 2\theta_k
\end{equation}
is applied as the final gradient step (multiply each accumulated component by
\texttt{2.0*theta[k]}).

% ============================================================================
\section{Algorithm Pseudocode}\label{sec:pseudo}
% ============================================================================

\begin{lstlisting}[caption={Truncated \& Renormalized ETAS fit.}]
Algorithm: Truncated & Renormalized ETAS Fit
=============================================
Input:  Catalogue (t,x,y,z,m), polygon, study window,
        eps_t, eps_s, eps_z, max_iter

1.  Initialize sqrt-parameters theta_0
2.  For iter = 1 to max_iter:
    a.  Compute T_max, R_max, G_norm, F_norm, H_norm
        and their total derivatives (renorm.py)
    b.  Build KDTree NeighborIndex and pre-prune parent lists
    c.  E-step: stochastic declustering -> background probs
    d.  M-step: DFP quasi-Newton optimization over theta
        - lambda_j and grad using KDTree-pruned parents
        - space-time integral Lambda_j and grad
        - Background rate update: mu = sum(rho_j)/N
3.  Return fitted parameters, Fisher information, etc.
\end{lstlisting}

% ============================================================================
\section{Implementation Notes}\label{sec:impl}
% ============================================================================

\subsection{Numerical Stability}

\begin{itemize}
\item \textbf{Base clamping in \texttt{\_safe\_pow}.} The kernels involve
      $(1+X)^{-\alpha}$ with $\alpha\notin\mathbb{Z}$, which can produce
      complex garbage when $1+X$ rounds slightly below $1$ in float64.
      \texttt{\_safe\_pow} clamps the base to $[1,\infty)$.
\item \textbf{Depth kernel \texttt{safe\_u}.} The logs $\ln u$ and $\ln(1-u)$
      are clamped at $10^{-12}$ to avoid $-\infty$ at the boundaries $u=0,1$.
\item \textbf{GPU safety.} The depth kernel uses \texttt{xp.where} instead of
      Python \texttt{max()} for compatibility with CuPy backends.
\end{itemize}

\subsection{Handling Invalid Parameters During Line Search}

The DFP line search can transiently push $p$ or $q$ below their lower bounds
($p\le 1$, $q\le 1$). Rather than raising exceptions, the cutoff and norm
functions return identity defaults ($\Tmax,\Rmax\to\infty$;
$\Gnorm,\Fnorm\to 1$; gradients $\to 0$). The optimizer detects the poor
likelihood value and backs off naturally.

\subsection{Testing Strategy}

\begin{center}
\begin{tabular}{@{}ll@{}}
\toprule
Test file & What it validates\\
\midrule
\texttt{test\_renorm.py}          & $\Gnorm,\Fnorm$ vs quadrature; total-derivative FD\\
\texttt{test\_neighbors.py}       & KDTree lookup vs brute-force masks\\
\texttt{test\_gradient.py}        & $\partial\lambda_j/\partial\theta$ vs central FD (2D, 3D, renorm)\\
\texttt{test\_3d\_regression.py}  & 3D = 2D when $h\equiv 1$; identity norms preserve output\\
\texttt{test\_pruning\_accuracy.py}& Pruned renormalized $\ell$ close to full; tighter $\eps$ shrinks error\\
\texttt{test\_roundtrip.py}       & Pipeline: catalog $\to$ fit $\to$ \texttt{par\_names} (2D \& 3D)\\
\bottomrule
\end{tabular}
\end{center}

% ============================================================================
\section{Sample Output}\label{sec:output}
% ============================================================================

The table below compares the baseline (un-pruned) ETAS fit against the
truncated \& renormalized model at two pruning thresholds on a synthetic
catalog of \textbf{288 events} (3 EM iterations, \texttt{mver=1}, Windows~11,
Python~3.13).

\begin{lstlisting}[caption={Baseline vs.\ truncated \& renormalized (288 events).}]
====================================================================================
  Catalog: 288 events (3 EM iterations)
  Model                              Time (s)         Log-lik   |Delta-LL|
  -------------------------------- ---------- --------------- ------------
  Baseline (no pruning)                 236.3       -237.5849            -
  Renormalized, eps = 1e-3              179.9       -238.6884       1.1035
  Renormalized, eps = 1e-5              178.9       -237.7117       0.1268
  Speedup Renormalized, eps = 1e-3       1.31x
  Speedup Renormalized, eps = 1e-5       1.32x

  Fitted parameters (mu, A, c, alpha, p, D, q, gamma):
  Baseline (no pruning)            [  1.02677,   0.46047,   2.41832,   0.58535,   3.23329,   0.05729,   4.15617,   0.06199]
  Renormalized, eps = 1e-3         [  1.02434,   0.45336,   4.91953,   0.56343,   5.15153,   0.06327,   4.45335,   0.05770]
  Renormalized, eps = 1e-5         [  1.01532,   0.45932,   2.56985,   0.58489,   3.34855,   0.05762,   4.17684,   0.06262]
====================================================================================
\end{lstlisting}

\paragraph{Key observations.}
\begin{enumerate}
\item \textbf{Renormalization is correct.} With $\eps=10^{-5}$ the
      log-likelihood differs from baseline by only $0.13$ ($\approx 0.05\%$
      of the baseline value) and the fitted parameters are nearly identical.
      The small remaining difference is the approximation error of the finite
      pruning threshold; it shrinks with smaller $\eps$ (Prop.~\ref{prop:limit}).
\item \textbf{Even small catalogs benefit.} For this 288-event catalogue the
      renormalized model achieves $1.3\times$ speedup by skipping the
      $O(N^2)$ pairwise mask construction. For larger catalogues
      ($N>10{,}000$) speedups of \textbf{10--50$\times$} are typical because
      KDTree lookups scale as $O(N\log N)$ while the brute-force mask scales
      as $O(N^2)$.
\item \textbf{Trading accuracy for speed.} At $\eps=10^{-3}$ the
      log-likelihood difference is $1.10$ ($\approx 0.46\%$). Users may choose
      $\eps$ to balance speed and precision.
\end{enumerate}

\paragraph{To reproduce:} run \texttt{python run\_comparison.py} from the
package root.

% ============================================================================
%  References — manual thebibliography (no .bib / no biber required).
%  Every \cite/\citep/\citet key in the text resolves to a \bibitem below.
%  Compiled citation style is author-year (natbib, round option).
%  Keys in use:  J19, GZZ15, O98
%  (ZOV02 is provided for completeness though not \cite'd in the body.)
% ============================================================================
\renewcommand{\refname}{References}
\begin{thebibliography}{99}
\addcontentsline{toc}{section}{References}

\bibitem[Guo et~al.(2015)]{GZZ15}
Guo, Y., Zhuang, J., and Zhou, S. (2015).
\newblock A hypocentral version of the space--time ETAS model.
\newblock \emph{Geophysical Journal International}, 203(1), 366--372.

\bibitem[Jalilian(2019)]{J19}
Jalilian, A. (2019).
\newblock ETAS: An R package for fitting the space-time ETAS model to
earthquake data.
\newblock \emph{Journal of Statistical Software}, 88, Code Snippet 1.

\bibitem[Ogata(1998)]{O98}
Ogata, Y. (1998).
\newblock Space-time point-process models for earthquake occurrences.
\newblock \emph{Annals of the Institute of Statistical Mathematics},
50(2), 379--402.

\bibitem[Zhuang et~al.(2002)]{ZOV02}
Zhuang, J., Ogata, Y., and Vere-Jones, D. (2002).
\newblock Stochastic declustering of space-time earthquake occurrences.
\newblock \emph{Journal of the American Statistical Association},
97(458), 369--380.

\end{thebibliography}

\end{document}
